\documentclass[11pt,letterpaper]{article}
% Reusable style. Build each .tex from its own directory using scripts/build.sh.
\usepackage[T1]{fontenc}
\usepackage[utf8]{inputenc}
\usepackage{lmodern}
\usepackage[margin=0.86in,headheight=15pt]{geometry}
\usepackage{amsmath,amssymb,amsthm,mathtools}
\usepackage{booktabs,longtable,tabularx,array,multicol}
\usepackage[table]{xcolor}
\usepackage{microtype}
\usepackage{enumitem}
\usepackage{fancyhdr}
\usepackage{titlesec}
\usepackage{hyperref}
\usepackage{listings}
\usepackage[most]{tcolorbox}
\usepackage{needspace}
\definecolor{aeNavy}{HTML}{172D43}
\definecolor{aeBlue}{HTML}{176B90}
\definecolor{aeTeal}{HTML}{148477}
\definecolor{aePale}{HTML}{F2F7F9}
\definecolor{aeGray}{HTML}{5C6771}
\hypersetup{colorlinks=true,linkcolor=aeBlue,urlcolor=aeBlue,pdfauthor={Algorithm Engineering Documentation Suite},pdfsubject={Independent educational materials based on the supplied TAOCP volume topics}}
\titleformat{\section}{\Large\bfseries\color{aeNavy}}{\thesection}{0.6em}{}
\titleformat{\subsection}{\large\bfseries\color{aeBlue}}{\thesubsection}{0.6em}{}
\titleformat{\subsubsection}{\normalsize\bfseries\color{aeNavy}}{\thesubsubsection}{0.6em}{}
\setlist[itemize]{leftmargin=*,topsep=3pt,itemsep=1pt}
\setlist[enumerate]{leftmargin=*,topsep=3pt,itemsep=3pt}
\setlength{\parindent}{0pt}\setlength{\parskip}{5pt}
\pagestyle{fancy}\fancyhf{}\fancyhead[L]{\footnotesize\color{aeGray}ALGORITHM ENGINEERING}\fancyhead[R]{\footnotesize\color{aeGray}DOCUMENTATION SUITE}\fancyfoot[L]{\footnotesize\color{aeGray}Independent study and engineering practice}\fancyfoot[R]{\footnotesize\thepage}
\renewcommand{\headrulewidth}{0.35pt}
\newtcolorbox{engineeringnote}[1][]{breakable,colback=aePale,colframe=aeBlue,boxrule=0.5pt,arc=1mm,fonttitle=\bfseries,title={Engineering note},#1}
\newtcolorbox{evidencebox}[1][]{breakable,colback=white,colframe=aeTeal,boxrule=0.6pt,arc=1mm,fonttitle=\bfseries,title={Evidence and deliverables},#1}
\lstdefinestyle{algo}{basicstyle=\ttfamily\footnotesize,breaklines=true,frame=single,rulecolor=\color{aeGray!40},backgroundcolor=\color{aePale},showstringspaces=false,tabsize=2,columns=fullflexible,keepspaces=true,keywordstyle=\color{aeBlue}\bfseries,commentstyle=\color{aeTeal}}
\lstset{style=algo}
\newcommand{\DocTitle}[3]{%
\hypersetup{pageanchor=false}\begin{titlepage}\vspace*{1.2cm}
{\color{aeTeal}\rule{\linewidth}{2.3pt}}\par\vspace{1.15cm}
{\fontsize{24}{30}\selectfont\bfseries\color{aeNavy}\raggedright #1\par}\vspace{0.45cm}
{\Large\color{aeBlue}#2\par}\vspace{1.3cm}
{\large\bfseries #3\par}\vspace{0.35cm}
{\color{aeGray}Edition 1.0 \quad | \quad 29 October 2026\par}\vfill
{\color{aeGray}\small Original educational material aligned to the chapter and section topics in the five supplied volumes of Donald E. Knuth's \emph{The Art of Computer Programming}. This document is not an authorized reproduction or substitute for the original books.}\par\vspace{0.4cm}
{\color{aeTeal}\rule{\linewidth}{0.75pt}}
\end{titlepage}\hypersetup{pageanchor=true}\pagenumbering{roman}}
\newcommand{\SourceNote}[1]{\par\smallskip{\footnotesize\color{aeGray}\textbf{Source locator:} #1}\par\smallskip}
\newcommand{\Outcome}[1]{\par\textbf{Competency target.} #1\par}
\newcommand{\Checkpoint}[1]{\par\textbf{Checkpoint.} #1\par}
% Intentionally simple printable writing lines: call \AnswerLines{number}.
\newcount\alcount
\newcommand{\AnswerLines}[1]{\par\alcount=0\loop\ifnum\alcount<#1\par\noindent\rule{\linewidth}{0.3pt}\vspace{0.33cm}\advance\alcount by1\repeat\par}
\newcommand{\FormField}[2][0.62\linewidth]{\noindent\textbf{#2}\hspace{0.65em}\rule{#1}{0.4pt}\par\vspace{0.18cm}}
\newcommand{\ReviewTable}{\begin{tabularx}{\linewidth}{@{}p{0.23\linewidth}X p{0.2\linewidth}@{}}\toprule Criterion & Recorded evidence & Rating (0--4)\\\midrule Correctness & & \\[1em] Proof quality & & \\[1em] Complexity & & \\[1em] Testing & & \\[1em] Reproducibility & & \\[1em] Communication & & \\[1em]\bottomrule\end{tabularx}}

\begin{document}
\DocTitle{Prerequisite-Based Learning Curriculum}{A 24-week project-driven route across all five volumes}{Engineering documentation}
\tableofcontents
\clearpage
\pagenumbering{arabic}
\section{Entry prerequisites and readiness checks}
The learner should be comfortable writing and testing small programs in Python and modern C++, using Git, reading loops and recursion, and doing elementary algebra. Readiness checks: derive a triangular sum, implement binary search, describe signed integer overflow in C++, and compile a C++ program with sanitizers. A deficiency triggers targeted study, not abandonment of the curriculum.
\section{Concept prerequisite graph}
\begin{verbatim}
Logic and induction -----> proof invariants -----> backtracking / SAT
        |                        |
Discrete counting ------> asymptotic analysis ------> sorting lower bounds
        |                        |
Arrays, pointers -------> storage structures ------> search and indexing
        |                        |
Integer arithmetic -----> number representation ---> PRNG and precision
        |                                                |
        +------------> combinatorial enumeration --------+
\end{verbatim}
Dependencies denote pedagogical readiness, not exclusive paths. Read the volume numbers as references rather than as fixed weekly homework.
\section{Sequenced learning plan}
\small
\begin{longtable}{@{}p{0.11\linewidth}p{0.24\linewidth}p{0.37\linewidth}p{0.20\linewidth}@{}}
\toprule Weeks & Primary source & Hands-on output & Exit condition\\\midrule\endfirsthead
\toprule Weeks & Primary source & Hands-on output & Exit condition\\\midrule\endhead
1--2 & Vol. 1, 1.1--1.2 & Contract and proof exercises & Valid invariant/variant\\
3--4 & Vol. 1, 1.3--1.4 & State-machine iterator & Tests + state trace\\
5--6 & Vol. 1, Ch. 2 & Ring buffer and tree traversal & Bound + edge tests\\
7--8 & Vol. 2, Ch. 3 & Seeded PRNG and shuffle study & Valid statistical caveats\\
9--10 & Vol. 2, Ch. 4 & GCD, base conversion, FP audit & Error report + proof\\
11--13 & Vol. 3, Ch. 5 & Insertion/merge sort comparison & Stability proof + data\\
14--15 & Vol. 3, Ch. 6 & Binary search and index prototype & Lookup contract + costs\\
16--18 & Vol. 4A, 7.1--7.2.1 & Bitsets, BDD sketch, enumeration & Exhaustive small checks\\
19--20 & Vol. 4B, 7.2.2 & Instrumented backtracking & Correct small-instance counts\\
21--22 & Vol. 4B, 7.2.2.1 & Exact-cover solver & Witness verified\\
23--24 & Vol. 4B, 7.2.2.2 & DPLL SAT solver and capstone & Validation + defense\\
\bottomrule\end{longtable}
\normalsize
\section{Weekly cadence}
\begin{enumerate}\item \textbf{Reading:} identify the main abstractions, prerequisites and at least one proof technique in the source sections.
\item \textbf{Derivation:} write a one-page contract, invariant and complexity calculation without consulting code.
\item \textbf{Implementation:} maintain independent Python and C++ versions under the same semantic specification.
\item \textbf{Verification:} run unit, property, negative and differential tests; log uncovered cases.
\item \textbf{Experimentation:} benchmark only after tests pass; keep raw data and machine manifest.
\item \textbf{Review:} complete the scorecard and enter outcomes into the progress tracker.
\end{enumerate}
\section{Three optional specialization paths}
\subsection{Systems and performance engineering}
Prioritize contiguous structures, cache-sensitive searching, external sorting, memory ownership, compiler optimization and profilers. Capstone: compare an in-memory index against disk-backed runs under bounded-memory assumptions.
\subsection{Numerical and computational engineering}
Prioritize number representation, GCD, floating-point error, polynomial evaluation, PRNG quality and reproducible experiments. Capstone: produce an error budget for a stochastic numerical workflow.
\subsection{Search and constraint engineering}
Prioritize counting, Boolean representation, permutations, exact cover, dynamic choice ordering and SAT. Capstone: encode a scheduling constraint as exact cover or CNF, validate it, and compare search-node counts.
\section{Mastery standard and portfolio}
A module is complete only when (1) the written specification is unambiguous, (2) correctness is justified or limitations are recorded, (3) Python and C++ outputs agree on a shared fixture, (4) time/space cost is derived, (5) experiment instructions reproduce outputs, and (6) a review scorecard reaches the agreed threshold. Retake failed gates with new adversarial inputs.
\end{document}
